% ---------------------------------------------------------------------------
% Seminar report: Real-Time Image-Based Lighting with the Split-Sum Approximation
% Advanced Computer Graphics, 2025/26 -- Topic 2.9
% ---------------------------------------------------------------------------

\documentclass{egpubl}
\JournalSubmission
\electronicVersion

\ifpdf \usepackage[pdftex]{graphicx} \pdfcompresslevel=9
\else  \usepackage[dvips]{graphicx} \fi

\PrintedOrElectronic

\usepackage{t1enc,dfadobe}
\usepackage{egweblnk}
\usepackage{cite}

\usepackage{amsmath}
\usepackage{booktabs}
\usepackage{listings}
\usepackage{xcolor}

\lstset{
    basicstyle=\ttfamily\scriptsize,
    breaklines=true,
    columns=fullflexible,
    showstringspaces=false,
    keywordstyle=\color{blue!70!black},
    commentstyle=\color{gray}
}

% ---------------------------------------------------------------------------
\title[Real-Time IBL with Split-Sum]%
      {Real-Time Image-Based Lighting with the Split-Sum Approximation}

\author[V. Poto\v{c}nik]
{\parbox{\textwidth}{\centering Vid Poto\v{c}nik$^{1}$}
 \\
 {\parbox{\textwidth}{\centering
   $^1$University of Ljubljana, Faculty of Computer and Information Science, Slovenia
 }}
}

% ---------------------------------------------------------------------------
\begin{document}
\maketitle

\begin{abstract}
\textit{We present a real-time rendering pipeline that uses image-based lighting (IBL)
to evaluate the contribution of an HDR environment map to physically based shading.
The pipeline precomputes a diffuse irradiance cubemap by cosine-weighted hemisphere
convolution, a specular prefiltered environment cubemap with mip-level encoded
roughness via GGX importance sampling, and a two-dimensional BRDF integration lookup
table using the split-sum approximation of Karis. At render time the diffuse and
specular IBL terms are evaluated with three texture fetches and a Schlick--Fresnel
modulation. We implement the pipeline in OpenGL~3.3 and analyse the trade-offs
between preprocessing time, memory cost, and rendering accuracy across three HDR
environments.}

\begin{CCSXML}
<ccs2012>
<concept>
<concept_id>10010147.10010371.10010372.10010374</concept_id>
<concept_desc>Computing methodologies~Rendering</concept_desc>
<concept_significance>500</concept_significance>
</concept>
<concept>
<concept_id>10010147.10010371.10010372.10010375</concept_id>
<concept_desc>Computing methodologies~Reflectance modeling</concept_desc>
<concept_significance>300</concept_significance>
</concept>
</ccs2012>
\end{CCSXML}

\ccsdesc[500]{Computing methodologies~Rendering}
\ccsdesc[300]{Computing methodologies~Reflectance modeling}

\printccsdesc
\end{abstract}

% ---------------------------------------------------------------------------
\section{Introduction}

Realistic surface appearance under arbitrary illumination depends strongly on the
surrounding environment. Image-based lighting (IBL), introduced by
Debevec~\cite{Debevec1998} and later codified as a practical
technique~\cite{Debevec2006}, uses a captured high-dynamic-range (HDR) panoramic
image as the sole light source, evaluating the rendering equation against the
panoramic radiance distribution directly. A naive evaluation requires Monte Carlo
integration per pixel, which is incompatible with real-time constraints.

Karis~\cite{Karis2013} proposed a \emph{split-sum} approximation that decomposes
the specular IBL integral into a precomputed prefiltered environment map and a 2D
scale--bias lookup table. Combined with the standard cosine-weighted irradiance
map for the diffuse term~\cite{Ramamoorthi2001}, this reduces IBL evaluation to
three constant-time texture fetches per pixel and is the foundation of the
real-time physically based rendering (PBR) pipelines used in modern game
engines~\cite{Burley2012}.

This work implements and analyses such a pipeline for the GGX microfacet
BRDF~\cite{Walter2007}. Section~\ref{sec:background} reviews the underlying theory.
Section~\ref{sec:pipeline} describes the four preprocessing passes and the per-frame
shading. Section~\ref{sec:impl} describes the implementation.
Section~\ref{sec:results} reports preprocessing times, rendering performance, and a
qualitative discussion of the accuracy/performance trade-offs.

% ---------------------------------------------------------------------------
\section{Background}\label{sec:background}

\subsection{The rendering equation under environmental lighting}

The outgoing radiance $L_o$ at a surface point in direction $\omega_o$ is
\begin{equation}
L_o(\omega_o) = \int_\Omega f_r(\omega_i, \omega_o)\,L_i(\omega_i)\,
                  (\mathbf{n}\cdot\omega_i)\,\mathrm{d}\omega_i,
\label{eq:rendering}
\end{equation}
where $L_i$ is the incoming radiance, $f_r$ is the bidirectional reflectance
distribution function (BRDF), $\mathbf{n}$ is the surface normal, and $\Omega$ is
the upper hemisphere. Under image-based lighting the scene is enclosed by an HDR
panoramic image at infinity; $L_i(\omega_i)$ is read directly from this image,
typically resampled into a cubemap to allow constant-time GPU access.

\subsection{Cook--Torrance microfacet BRDF}

We adopt the Cook--Torrance microfacet BRDF~\cite{CookTorrance1982} with the GGX
normal distribution function~\cite{Walter2007}, the height-correlated Smith
geometry term~\cite{Heitz2014}, and the Schlick approximation of the Fresnel term:
\begin{equation}
f_r =
\underbrace{\frac{(1-F)(1-m)\,\mathbf{c}}{\pi}}_{\text{Lambertian}}
\;+\;
\underbrace{\frac{D(\mathbf{h})\,F(\omega_o,\mathbf{h})\,G(\omega_i,\omega_o)}
                 {4(\mathbf{n}\cdot\omega_i)(\mathbf{n}\cdot\omega_o)}}_{\text{specular}}.
\label{eq:brdf}
\end{equation}
Here $\mathbf{h}$ is the half-vector, $m \in [0,1]$ is the metallic factor,
$\mathbf{c}$ the base colour (albedo), and $F_0$ the reflectance at normal
incidence: $F_0 = \text{mix}(0.04, \mathbf{c}, m)$. The factor $(1-F)$
enforces energy conservation between the diffuse and specular lobes, and the
factor $(1-m)$ removes the diffuse contribution for metallic surfaces.

\subsection{The split-sum approximation}

Equation~\eqref{eq:rendering} cannot be evaluated analytically against an HDR
environment, but it admits a useful factorisation for the two lobes of
\eqref{eq:brdf}.

\paragraph*{Diffuse term.} For a Lambertian surface, $f_r$ is constant and the
integral reduces to the cosine-weighted average of incoming radiance,
\begin{equation}
E(\mathbf{n}) = \int_\Omega L_i(\omega_i)\,(\mathbf{n}\cdot\omega_i)\,\mathrm{d}\omega_i.
\end{equation}
$E$ depends only on the surface normal and can be precomputed once into a cubemap
indexed by $\mathbf{n}$, the \emph{irradiance map}. The angular bandlimit imposed
by the cosine kernel is low, so a $32\!\times\!32$ resolution per face is
sufficient.

\paragraph*{Specular term.} The specular integrand does not separate cleanly,
because $f_r^{\text{spec}}$ depends jointly on $\omega_i$ and $\omega_o$.
Karis~\cite{Karis2013} approximates it as a product of two simpler integrals:
\begin{align}
&\int_\Omega f_r^{\text{spec}}(\omega_i,\omega_o)\,L_i(\omega_i)\,(\mathbf{n}\cdot\omega_i)\,\mathrm{d}\omega_i \nonumber \\
&\quad \approx
\underbrace{\Bigg(\frac{\sum_{k=1}^{N}\!L_i(\omega_k)(\mathbf{n}\cdot\omega_k)}
                       {\sum_{k=1}^{N}\!(\mathbf{n}\cdot\omega_k)}\Bigg)}_{\text{(I) prefiltered environment}}
\;\cdot\;
\underbrace{\int_\Omega f_r^{\text{spec}}(\omega_i,\omega_o)\,(\mathbf{n}\cdot\omega_i)\,\mathrm{d}\omega_i}_{\text{(II) BRDF integral}}\!.
\label{eq:split}
\end{align}
Factor~(I) depends on the reflection direction $\mathbf{r} = \text{reflect}(-\omega_o,\mathbf{n})$
and roughness $\alpha$. It is precomputed into a roughness-indexed sequence of
cubemaps, encoded as mip levels of a single cubemap texture: mip~0 holds the
mirror-reflection case ($\alpha=0$) and mip~4 holds the fully rough case
($\alpha=1$). The samples $\omega_k$ are drawn by importance-sampling the GGX
NDF using a Hammersley quasi-random sequence.

Factor~(II) depends only on $(\mathbf{n}\cdot\omega_o, \alpha)$ and admits a
further algebraic simplification by isolating the Fresnel term. Substituting
$F(\omega_o,\mathbf{h}) = F_0 + (1-F_0)(1-\omega_o\!\cdot\!\mathbf{h})^5$ yields
\begin{equation}
\int_\Omega f_r^{\text{spec}}(\mathbf{n}\cdot\omega_i)\,\mathrm{d}\omega_i
\;=\;
F_0\,A(\mathbf{n}\!\cdot\!\omega_o,\alpha) + B(\mathbf{n}\!\cdot\!\omega_o,\alpha),
\label{eq:lut}
\end{equation}
where $A$ and $B$ are scalar functions of two arguments. They are precomputed
into a single 2D lookup table, the \emph{BRDF LUT}, with the two channels of an
\texttt{RG16F} texture storing $A$ and $B$ respectively.

At render time the specular IBL evaluation collapses to two texture fetches and
a multiply--add, independent of the environment complexity.

% ---------------------------------------------------------------------------
\section{Pipeline}\label{sec:pipeline}

\begin{figure}[t]
\centering
% \includegraphics[width=0.95\linewidth]{figures/pipeline.pdf}
\fbox{\parbox{0.92\linewidth}{\footnotesize\centering\vspace{1ex}
\texttt{HDR file} $\rightarrow$ \texttt{equirect texture} $\rightarrow$
\texttt{environment cubemap (512$^2$)} \\[0.5ex]
$\downarrow\qquad\qquad\quad\downarrow\qquad\qquad\quad\downarrow$ \\[0.5ex]
\texttt{irradiance (32$^2$)}\quad
\texttt{prefilter (128$^2$+mips)}\quad
\texttt{BRDF LUT (512$^2$)} \\[0.5ex]
$\downarrow$ \\[0.5ex]
\texttt{PBR fragment shader} $\rightarrow$ \texttt{tonemap} $\rightarrow$ \texttt{framebuffer}
\vspace{1ex}}}
\caption{Pipeline overview. The HDR equirectangular image is converted into a
cubemap, from which three precomputed structures are derived; the rendering pass
combines them with a Cook--Torrance Fresnel response.}
\label{fig:pipeline}
\end{figure}

The pipeline (Fig.~\ref{fig:pipeline}) consists of four preprocessing passes
executed once when an HDR environment is loaded, and a single fragment shading
pass per frame.

\subsection{Equirectangular--to--cubemap conversion}

The HDR file is loaded as a 2D texture in equirectangular projection (longitude
$\times$ latitude). A unit cube is rendered six times into the faces of a $512^2$
cubemap; the fragment shader samples the equirectangular texture using the texel
direction $\mathbf{d} = \widehat{\text{localPos}}$ converted to spherical UVs:
\begin{equation}
u = \tfrac{1}{2\pi}\arctan(\mathbf{d}_z, \mathbf{d}_x) + 0.5,\quad
v = \tfrac{1}{\pi}\arcsin(\mathbf{d}_y) + 0.5.
\end{equation}
Storage is \texttt{GL\_RGB16F} to retain the HDR range across the full pipeline.

\subsection{Diffuse irradiance convolution}

For each output texel of a $32^2$ cubemap face, the upper hemisphere around the
texel's direction is integrated by a uniform Riemann sum on $(\phi,\theta)$ with
a step of $0.025$\,rad ($\approx 1{,}300$ samples per texel). The integrand is
$L_i(\omega)\sin\theta\cos\theta$; the multiplicative $\pi$ from the diffuse
normalisation is applied at the end. The low output resolution suffices because
the cosine kernel acts as a low-pass filter that removes high-frequency content.

\subsection{GGX prefiltered environment map}

The prefiltered cubemap is allocated at $128^2$ with five mip levels and
\texttt{GL\_LINEAR\_MIPMAP\_LINEAR} filtering. For each mip level $i\in\{0,\dots,4\}$
the roughness is set to $\alpha = i/4$. The fragment shader, given the output
direction $\mathbf{n}$ (assumed equal to $\mathbf{v}$ and $\mathbf{r}$ following
Karis), draws $N$ samples (default $N=1024$) from a Hammersley
sequence~\cite{Hammersley1964}, transforms them into half-vectors $\mathbf{h}$
through the GGX inverse-CDF, computes $\mathbf{L} = \text{reflect}(-\mathbf{v},\mathbf{h})$,
and accumulates the environment radiance weighted by $(\mathbf{n}\cdot\mathbf{L})$.

\begin{lstlisting}[caption={GGX importance sampling for the prefilter pass.},
                   label={lst:prefilter}]
vec3 importanceSampleGGX(vec2 Xi, vec3 N, float roughness) {
    float a   = roughness * roughness;
    float phi = 2.0 * PI * Xi.x;
    float ct  = sqrt((1.0 - Xi.y) / (1.0 + (a*a - 1.0)*Xi.y));
    float st  = sqrt(1.0 - ct*ct);
    vec3  H   = vec3(cos(phi)*st, sin(phi)*st, ct);
    // ...transform from tangent space into world space using N...
    return normalize(tangent*H.x + bitangent*H.y + N*H.z);
}
\end{lstlisting}

\subsection{BRDF lookup table}

A $512\!\times\!512$ \texttt{RG16F} texture is rendered once at startup with a
single full-screen quad. The fragment shader evaluates equation~\eqref{eq:lut}
by importance-sampling the GGX NDF with 1024 Hammersley samples and accumulating
$(1{-}F_c)\,G_{\text{vis}}$ and $F_c\,G_{\text{vis}}$ separately into the two
output channels, where $F_c = (1{-}\mathbf{v}\!\cdot\!\mathbf{h})^5$ and
$G_{\text{vis}}$ is the visibility-weighted Smith geometry term. The result is
material-independent and need not be regenerated when the environment or the
material changes.

\subsection{Rendering pass}

Per pixel the fragment shader (Listing~\ref{lst:pbr}) computes the diffuse
contribution with one cubemap fetch, the specular contribution with two (the
prefiltered map sampled at \textsc{lod}$=\alpha\cdot4$ and the BRDF LUT at
$(\mathbf{n}\!\cdot\!\mathbf{v}, \alpha)$), and combines them with a
roughness-modulated Schlick Fresnel response~\cite{Lazarov2013} that avoids
edge darkening at high roughness.

\begin{lstlisting}[caption={Core of the per-pixel IBL evaluation.},
                   label={lst:pbr}]
vec3 F  = fresnelSchlickRoughness(NdotV, F0, roughness);
vec3 kD = (1.0 - F) * (1.0 - metallic);

vec3 irradiance = texture(irradianceMap, N).rgb;
vec3 diffuse    = kD * irradiance * albedo;

vec3 prefiltered = textureLod(prefilterMap, R, roughness * 4.0).rgb;
vec2 brdf        = texture(brdfLUT, vec2(NdotV, roughness)).rg;
vec3 specular    = prefiltered * (F * brdf.x + brdf.y);

color = (diffuse + specular) * exposure;
color = tonemap(color);
\end{lstlisting}

The final colour passes through a tonemapping operator (Reinhard, ACES filmic,
or Uncharted~2~\cite{Hable2010}) and gamma correction before being written to
the framebuffer.

% ---------------------------------------------------------------------------
\section{Implementation}\label{sec:impl}

The application is implemented in C++17 against OpenGL~3.3 core profile, with
GLFW for window/context management, GLAD as the function loader, GLM for linear
algebra, \texttt{stb\_image} for HDR decoding, and Dear ImGui for the parameter
UI. All map dimensions and the prefilter sample count are exposed through ImGui
(Fig.~\ref{fig:ui}); modifying any of them triggers regeneration of the affected
map without restarting the application. Three primitive meshes (sphere, cube,
torus) and a $7\!\times\!7$ roughness--metallic grid mode are provided for visual
inspection. The source tree is approximately 1.5\,kLOC of C++ and 350~LOC of GLSL,
and builds to a single executable through CMake.

\begin{figure}[t]
\centering
% \includegraphics[width=0.55\linewidth]{figures/ui.png}
\fbox{\parbox{0.55\linewidth}{\footnotesize\centering\vspace{0.5ex}
[screenshot of the ImGui panel]
\vspace{0.5ex}}}
\caption{ImGui control panel. The user adjusts material parameters, exposure,
the tonemapping operator, scene composition (single object or material grid),
the active HDR environment, and the four IBL preprocessing resolutions.}
\label{fig:ui}
\end{figure}

% ---------------------------------------------------------------------------
\section{Results}\label{sec:results}

\subsection{Qualitative evaluation}

Figure~\ref{fig:grid} shows the $7\!\times\!7$ material grid under
\texttt{abandoned\_parking} (rows: metallic $0\!\rightarrow\!1$;
columns: roughness $0.05\!\rightarrow\!1$). Dielectric materials (top row)
exhibit the expected Fresnel rim brightening and roughness-dependent reflection
blur. Pure metals (bottom row) show a coloured specular response controlled by
the albedo, with no diffuse contribution. Intermediate rows correctly interpolate
between the two regimes following the metallic workflow.

Figure~\ref{fig:env} contrasts a single mirror sphere under three HDR
environments. The pipeline correctly reproduces the overcast soft lighting of
\texttt{abandoned\_parking}, the studio key-and-fill lighting of
\texttt{studio\_small\_03}, and the strong sun-and-sky contrast of
\texttt{kloppenheim\_06}. Because $L_i$ is
stored at \texttt{GL\_RGB16F}, the sun disk in \texttt{kloppenheim\_06} retains
its high dynamic range and produces a clipped highlight after tonemapping,
matching the behaviour of physically captured reflections.

\begin{figure}[t]
\centering
% \includegraphics[width=\linewidth]{figures/grid_parking.jpg}
\fbox{\parbox{0.95\linewidth}{\footnotesize\centering\vspace{4ex}
[7$\times$7 material grid under \texttt{abandoned\_parking}]
\vspace{4ex}}}
\caption{Material grid. Rows top--to--bottom: increasing metallic.
Columns left--to--right: increasing roughness. A single HDR environment
(\texttt{abandoned\_parking}) is used for all 49 spheres.}
\label{fig:grid}
\end{figure}

\begin{figure}[t]
\centering
% \includegraphics[width=\linewidth]{figures/env_comparison.jpg}
\fbox{\parbox{0.95\linewidth}{\footnotesize\centering\vspace{2ex}
[mirror sphere under three HDR environments]
\vspace{2ex}}}
\caption{Mirror sphere ($\alpha=0.05$, $m=1$, white albedo) under three HDR
environments: \texttt{abandoned\_parking}, \texttt{studio\_small\_03},
and \texttt{kloppenheim\_06} (left to right).}
\label{fig:env}
\end{figure}

\subsection{Preprocessing performance}

Wall-clock preprocessing times are measured around a \texttt{glFinish()} barrier
to capture the actual GPU execution cost. Measurements are taken on an Apple
Silicon Mac with the OpenGL~4.1 Metal driver (Metal driver version 90.5).
Table~\ref{tab:sweep} sweeps the four independent parameters of the pipeline to
isolate the cost of each pass.

\begin{table*}[t]
\centering
\caption{Preprocessing time (ms) at five operating points on an Apple Silicon
Mac, single HDR environment, measured with \texttt{glFinish()} barriers. Row~1
includes a one-time shader-compilation overhead on first use; rows~2--5 are warm.}
\label{tab:sweep}
\small
\begin{tabular}{lcccc|rrrr|r}
\toprule
\# & Cubemap & Irradiance & Prefilter & Samples & Equirect & Irradiance & Prefilter & BRDF LUT & Total \\
\midrule
1 (cold) & $512^2$  & $32^2$ & $128^2$ & 1024 & 59.8 & 175.9 & 43.3 & 21.2 & 300.1 \\
2        & $1024^2$ & $32^2$ & $128^2$ & 1024 & 29.4 &  43.5 & 22.0 &  8.1 & 102.9 \\
3        & $1024^2$ & $64^2$ & $128^2$ & 1024 & 26.2 &  65.4 & 21.8 &  7.9 & 121.4 \\
4        & $1024^2$ & $64^2$ & $256^2$ & 1024 & 57.9 &  63.2 & 33.5 &  7.9 & 162.5 \\
5        & $1024^2$ & $64^2$ & $256^2$ & 2048 & 23.7 &  67.2 & 63.2 &  7.9 & 162.0 \\
\bottomrule
\end{tabular}
\end{table*}

Four observations follow directly from the sweep. First, doubling the prefilter
output resolution from $128^2$ to $256^2$ (rows~3$\rightarrow$4) increases the
prefilter pass from 21.8 to 33.5\,ms (1.5$\times$ rather than 4$\times$, because
the highest-resolution mip is a small fraction of the total texel count and the
lower mips remain at their original size). Second, doubling the prefilter sample
count from 1024 to 2048 (rows~4$\rightarrow$5) increases the prefilter pass from
33.5 to 63.2\,ms (1.9$\times$, close to the expected linear scaling). Third,
doubling the irradiance resolution from $32^2$ to $64^2$ (rows~2$\rightarrow$3)
takes the irradiance pass from 43.5 to 65.4\,ms (1.5$\times$); the irradiance
shader's inner Riemann loop is the same per texel, so the increase reflects the
$4\times$ growth in texel count partly absorbed by the GPU's spare capacity at
small kernel sizes. Fourth, the BRDF LUT cost is constant at $\sim$8\,ms
regardless of any other parameter, as expected since it is material- and
environment-independent.

Even at the most expensive operating point the total preprocessing cost remains
well under one second, which is acceptable for an interactive environment
switch.

\subsection{Rendering performance}

The PBR fragment shader performs a constant amount of work per pixel: three
texture fetches, two normalisations, one half-vector evaluation, and the
Cook--Torrance arithmetic. At $1280\!\times\!720$ on the same machine the
application runs at a sustained 120\,FPS in both the single-sphere mode and the
$7\!\times\!7$ grid mode (49 spheres, $\sim 0.7$\,M triangles after meshing at
$64\!\times\!64$ segments per sphere). This rate is the refresh-locked ceiling
of the ProMotion display and represents an upper bound rather than a measure of
the shader's intrinsic cost; nevertheless it confirms that the per-pixel
evaluation of equation~\eqref{eq:split} is well within the real-time budget
even when the rest of the geometry pipeline is moderately loaded.

\subsection{Trade-offs}

The split-sum factorisation introduces a small bias that is most visible at
grazing angles on rough surfaces; the assumption
$\mathbf{n}\!=\!\mathbf{v}\!=\!\mathbf{r}$ during the prefilter pass loses the
elongation of the GGX specular lobe in oblique reflections. In practice this
error is masked by the high entropy of natural HDR content and is generally not
perceptible against the more dominant effects of tonemapping and exposure.

The free parameters of the pipeline interact as follows:
\begin{itemize}
\setlength{\itemsep}{0pt}
\item The \emph{cubemap resolution} controls the angular fidelity of the source
      environment and bounds the maximum quality reachable by the prefilter pass.
\item The \emph{irradiance resolution} can be very low (16--64) without any
      perceptible loss, because the cosine kernel is strongly low-pass.
\item The \emph{prefilter base resolution} and \emph{mip count} control the
      angular precision of mirror-like reflections; the lowest-roughness mip
      contains the highest-frequency content and benefits most from increased
      resolution.
\item The \emph{importance sample count} controls residual high-frequency noise
      at intermediate-to-high roughness; for $\alpha \geq 0.5$ the angular
      density of useful samples drops sharply, and 2048 samples remain a useful
      upper bound on production engines.
\end{itemize}
Of these, only the prefilter cost grows linearly with the sample count and
quadratically with the base resolution; the others are negligible.

% ---------------------------------------------------------------------------
\section{Limitations and Future Work}

The implementation evaluates a single bounce of environmental light. It does not
model parallax inside the environment (the HDR map is assumed to be at infinity),
local shadowing of the IBL contribution (which would require an ambient occlusion
or screen-space visibility pass), or interreflection between scene objects. The diffuse irradiance is integrated by a uniform
$(\phi,\theta)$ Riemann sum rather than projected onto spherical
harmonics~\cite{Ramamoorthi2001}; the latter would compress the storage from a
cubemap to a 9-coefficient vector at no perceptual cost for the diffuse term.

The Karis prefilter assumption $\mathbf{n}=\mathbf{v}$ is the principal source
of approximation error in the specular term. It can be partially mitigated by
the multi-scattering formulation of Fdez-Ag{\"u}era~\cite{Fdez2019}, which
recovers the energy lost by the single-scattering Smith geometry term and is
critical for high-roughness metals.

% ---------------------------------------------------------------------------
\section{Conclusion}

We described a complete real-time IBL pipeline based on the split-sum
approximation, validated its behaviour qualitatively across three HDR
environments with different lighting characteristics, and quantified its
preprocessing and rendering cost. On an Apple Silicon Mac the technique attains
the vsync-locked frame rate of 120\,FPS at $1280\!\times\!720$ with a one-time
preprocessing budget of approximately 100--160\,ms per environment switch in
warm state, and exposes a small set of intuitive parameters that allow the user
to trade preprocessing time against angular fidelity. The source code, build instructions,
and the data used to produce this report are available under the MIT licence.

% ---------------------------------------------------------------------------
\bibliographystyle{eg-alpha-doi}
\bibliography{ibl-seminar}

\end{document}
